\centerline{\bf \Huge Турбо тренировка VI}

\begin{flushright}
        \scriptsize{\it мечта такая забавная штука,
                    \\ она может стать испытанием для смелости человека,
                    \\ а может стать местом, где он будет трусливо прятаться.
                    \\ Гриффит.}
\end{flushright}

\problem{1}
Дан выпуклый четырёхугольник периметра $10^{100}$, у которого длины всех сторон - натуральные числа, а сумма длин любых трёх сторон делится на длину оставшейся четвёртой стороны. Докажите, что этот четырёхугольник - ромб.

\problem{2}
Числа $b>0$ и $a$ таковы, что квадратный трехчлен $x^2+a x+b$ имеет два различных корня, ровно один из которых лежит на отрезке $[-1 ; 1]$. Докажите, что ровно один из этих корней лежит в интервале $(-b ; b)$.


\problem{3}
В стране некоторые пары городов соединены односторонними прямыми авиарейсами (между любыми двумя городами есть не более одного рейса). Скажем, что город $A$ доступен для города $B$, если из $B$ можно долететь в $A$, возможно, с пересадками. Известно, что для любых двух городов $P$ и $Q$ существует город $R$, для которого и $P$, и $Q$ доступны. Докажите, что существует город, для которого доступны все города страны. (Считается, что из города можно долететь до него самого.)

\problem{4}
Высоты остроугольного треугольника $A B C$, в котором $A B <A C$, пересекаются в точке $H$, а $O$ --- центр описанной около него окружности $\Omega$. Отрезок $O H$ пересекает описанную около треугольника $B H C$ окружность в точке $X$, отличной от $O$ и $H$. Окружность, описанная около треугольника $A O X$, пересекает меньшую дугу $A B$ окружности $\Omega$ в точке $Y$. Докажите, что прямая $X Y$ делит отрезок $B C$ пополам.